% Part: proof-theory
% Chapter: sequent-calculus
% Section: rules-proofs

\documentclass[../../../include/open-logic-section]{subfiles}

\begin{document}

\olfileid{pt}{seq}{rp}

\olsection{Rules and \usetoken{P}{proof}}

We begin by giving some definitions and fixing some terminology.

Consider, for example, the two rules in \Log{G3c} that allow !!{proving}
sequents containing $!A \land !B$ from sequents containing $!A$
and~$!B$:
\[
\Axiom$ !A, !B, \Gamma \fCenter \Delta$
\RightLabel{\LeftR{\land}}
\UnaryInf$ !A \land !B, \Gamma \fCenter \Delta$
\DisplayProof
\qquad
\Axiom$\Gamma \fCenter \Delta, !A$
\Axiom$ \Gamma \fCenter \Delta, !B$
\RightLabel{\RightR{\land}}
\BinaryInf$ \Gamma \fCenter \Delta, !A \land !B $
\DisplayProof
\]
The sequents on top of the line are called the \emph{premises} and the
sequent below the line the \emph{conclusion}. The !!{formula}s in
$\Gamma$ and $\Delta$ in each rule are called the \emph{context} or
the \emph{side !!{formula}s}. The !!{formula} in the conclusion that
isn't part of the context is called the \emph{principal} !!{formula},
and the !!{formula}s in the premises not in the context the
\emph{active} (or \emph{auxiliary}) !!{formula}s.

The quantifier rules are slightly more complicated. E.g., the rules in
\Log{G1c} for $\lforall$ are:
\[\Axiom$ !A(t), \Gamma \fCenter \Delta$
\RightLabel{\LeftR{\lforall}}
\UnaryInf$ \lforall[x][!A(x)],\Gamma \fCenter \Delta$
\DisplayProof
\qquad
\Axiom$ \Gamma \fCenter \Delta, !A(a) $
\RightLabel{\RightR{\lforall}}
\UnaryInf$ \Gamma \fCenter \Delta, \lforall[x][!A(x)]$
\DisplayProof
\]
In the \LeftR{\lforall} rule, the active !!{formula} is~$!A(t)$ where
$t$~is a closed term, i.e., a term without variables. Such an
inference is correct---that is, it matches the schematic rule
\LeftR{\lforall}---if there is some !!{formula}~$!A$, some
variable~$x$, and some closed term~$t$ such that the principal
!!{formula} in the conclusion is $\lforall[x][!A]$ and the active
!!{formula} in the premise is~$\Subst{!A}{t}{x}$. The
\RightR{\lforall} rule works the same, except that instead of any
term~$t$ the active !!{formula} in the premise must
be~$\Subst{!A}{a}{x}$, where $a$ is !!a{constant}. This constant is
called the \emph{eigenvariable} of the \RightR{\lforall} inference.
The inference is only correct if the lower sequent does not contain~$a$,
that is, if $a$ does not occur in $\Gamma$, $\Delta$, or~$!A$. This is
called the \emph{eigenvariable condition}.

Sequent systems like \Log{G1c} and \Log{G3c} consists of a collection
of such schematically specified rules together with some (also
schematically specified) sequents that count as axioms.
!!^{proof}s in such a system are finite trees of sequents that
are inductively generated from axioms using the rules of inference.
Each leaf is an axiom, and every other sequent follows from the
sequents immediately above it by one of the inference rules of the
system.

\begin{defn}\ollabel{defn:proof}
\emph{!!^{proof}s} in a sequent system are finite trees of
sequents inductively as follows:
\begin{enumerate}
  \item\ollabel{defn:prf:axiom} Every axiom is !!a{proof}.
  \item\ollabel{defn:prf:one-prem} If $\pi_1$ is !!a{proof} ending in
  a sequent~$S_1$, and $S$ is a sequent that follows from~$S_1$ by a
  one-premise rule~$R$, then
  \[
  \AxiomC{}
  \RightLabel{$\pi_1$}
  \DeduceC{$S_1$}
  \RightLabel{$R$}
  \UnaryInfC{$S$}
  \DisplayProof
  \]
  is !!a{proof}.
  \item\ollabel{defn:prf:two-prem} If $\pi_1$ and $\pi_2$ are !!{proof}s ending in sequents~$S_1$ and~$S_2$, and $S$ is a
  sequent that follows from~$S_1$ and$S_2$ by a two-premise rule~$R$, then
  \[
  \AxiomC{}
  \RightLabel{$\pi_1$}
  \DeduceC{$S_1$}
  \AxiomC{}
  \RightLabel{$\pi_2$}
  \DeduceC{$S_2$}
  \RightLabel{$R$}
  \BinaryInfC{$S$}
  \DisplayProof
  \]
  is !!a{proof}.
\end{enumerate}
\end{defn}

\Olref{defn:proof} defines !!{proof}s as trees of sequents. We think
of these trees as growing ``upwards.'' The topmost (leaf) nodes of
this tree are axioms and called \emph{initial sequents}. The
bottom-most sequent is called the \emph{end-sequent}. If $\pi$ is
!!a{proof} with end-sequent~$S$ we also write $\pi \Proves S$. If
there is !!a{proof} of a sequent~$S$ we write $\Proves S$, possibly
indicating the system in which the proof exists on the left of the
$\Proves$~symbol, e.g., $\Log{G1c} \Proves !A, !B \Sequent !A \land
!B$. Every sequent in !!a{proof}, with the exception of initial
sequents, is the \emph{conclusion} of an inference by some rule~$R$.
The one or two sequents standing immediately above a sequent in
!!a{proof} (i.e., the parents of the node) are called the
\emph{premises} of the inference.

The !!{proof}s used in the inductive construction of
!!a{proof} are called its \emph{sub-!!{proof}s}. The
sub-!!{proof}s ending in the premises of an inference are called
the \emph{immediate sub-!!{proof}s} of the inference.

It is often necessary to measure the complexity of !!{proof}s as
a natural number. We could simply count the number of sequents in
!!a{proof}, but often a more useful measure is the !!{height} of
!!a{proof}: it is the maximal number of sequents between the
end-sequent and an initial sequent. We can also give an inductive
definition.

\begin{defn}
The \emph{!!{height}}~$\pheight{\pi}$ of !!a{proof}~$\pi$ is
inductively defined as follows.
\begin{enumerate}
  \item If $\pi$ consists of a single axiom, $\pheight{\pi} = 0$.
  \item If $\pi$ ends in an inference with one premise with immediate
  sub-!!{proof}~$\pi_1$, then $\pheight{\pi} = \pheight{\pi_1} + 1$.
  \item If $\pi$ ends in an inference with two premises with immediate
  sub-!!{proof}s~$\pi_1$ and~$\pi_2$, then $\pheight{\pi} =
  \max(\pheight{\pi_1}, \pheight{\pi_2}) + 1$.
\end{enumerate}
If a sequent~$S$ has !!a{proof} of height $\le n$, we write $\Proves[n] S$.
\end{defn}

\begin{ex}
In \Log{G3c}, any sequent of the form $!C, \Gamma \Sequent \Delta, !C$
is an axiom, where $!C$ must be atomic. Suppose $!D$ and $!E$ are
atomic sentences. Then $!D, !E \Sequent !D$ and $!D, !E \Sequent !E$
are instances of the axiom $!C, \Gamma \Sequent \Delta, !C$. E.g., if
we take $!C \ident !E$, $\Gamma = \{!D\}$ and $\Delta = \emptyset$ in
$!C, \Gamma \Sequent \Delta, !C$ we get exactly the sequent $!D, !E
\Sequent !E$. (Remember that we are dealing with multi\emph{sets} and
so $!D, !E \Sequent !E$ and $!E, !D \Sequent !D$ are the same
sequent.)

Since $!D, !E \Sequent !D$ is an instance of an axiom of~\Log{G3c}, it
(by itself) counts as !!a{proof} in~\Log{G3c}, and so does $!D, !E
\Sequent !E$ by \olref{defn:proof}\olref{defn:prf:axiom}. Let's
call them $\pi_1$ and~$\pi_2$. By \olref{defn:prf:two-prem}, we can
take these two !!{proof}s and generate a new one ($\pi_3$) using the
two-premise rule~$\RightR{\land}$:
\[
\Axiom$!D, !E \fCenter !D$
\Axiom$!D, !E \fCenter !E$
\RightLabel{\RightR{\land}}
\BinaryInf$!D, !E  \fCenter !D \land !E $
\DisplayProof
\]
From $\pi_3$ in turn we obtain the !!{proof}~$\pi_4$
\[
\Axiom$!D, !E \fCenter !D$
\Axiom$!D, !E \fCenter !E$
\RightLabel{\RightR{\land}}
\insertBetweenHyps{\hspace{5ex}}
\BinaryInf$!D, !E  \fCenter !D \land !E $
\LeftSubproofLabel{$\pi_3$}
\RightLabel{\LeftR{\land}}
\UnaryInf$!D \land !E  \fCenter !D \land !E$
\RightSubproofLabel{$\pi_4$}
\DisplayProof
\]
using the one-premise rule~$\LeftR{\land}$ (using
\olref{defn:proof}\olref{defn:prf:one-prem}). The immediate
sub-!!{proof} of the last inference in~$\pi_4$ is~$\pi_3$. The
immediate sub-!!{proof}s of the $\RightR{\land}$ inference are $\pi_1$
and~$\pi_2$. The sub-!!{proof}s of $\pi_4$ are all of $\pi_1$,
$\pi_2$, $\pi_3$, and $\pi_4$ itself. We have $\pheight{\pi_1} =
\pheight{\pi_2} = 0$, $\pheight{\pi_3} = 1$ and $\pheight{\pi_4} = 2$.
We have shown that $\Log{G3c} \Proves[4] !D \land !E  \fCenter !D
\land !E$ for any atomic $!D$ and~$!E$.
\end{ex}

\end{document}